% Generated draft: requires mathematical and source review.
\documentclass[11pt]{article}
\usepackage[margin=25mm]{geometry}
\usepackage{fontspec}
\setmainfont{Latin Modern Roman}
\usepackage{amsmath,amssymb,mathtools,unicode-math}
\setmathfont{Latin Modern Math}
\usepackage{enumitem,needspace,xurl,hyperref}
\setlength{\parindent}{0pt}
\setlength{\parskip}{5pt}
\setlength{\emergencystretch}{3em}
\newcommand{\N}{\mathbb{N}}
\newcommand{\Z}{\mathbb{Z}}
\newcommand{\Q}{\mathbb{Q}}
\newcommand{\R}{\mathbb{R}}
\newcommand{\C}{\mathbb{C}}
\newcommand{\F}{\mathbb{F}}
\newcommand{\A}{\mathbb{A}}
\newcommand{\PP}{\mathbb{P}}
\newcommand{\E}{\mathbb{E}}
\newcommand{\eps}{\varepsilon}
\newcommand{\dd}{\,\mathrm{d}}
\begin{document}
\section*{UnsolvedMath ID: 700708}
\textbf{Source problem code:} catalogue entry
\subsection{Definitions and mathematical statement}
Let \(F(x,y)\) be the free group on two generators. A word \(w\in F(x,y)\) defines an evaluation \(w(g,h)\in G\) for every group \(G\) and every \(g,h\in G\). Write \(w(G)=1\) when \(w(g,h)=1_G\) for all \(g,h\in G\). A Sylow \(p\)-subgroup of a finite group \(G\) is a subgroup of order equal to the largest power of \(p\) dividing \(|G|\); this includes the trivial subgroup when \(p\nmid |G|\).

The source states that for every prime \(p\) there exists a sequence \((a_{p,k})_{k\geq 1}\) of words in \(F(x,y)\), depending on \(p\) but independent of \(G\), such that every finite group \(G\) satisfies
\[
\text{every Sylow \(p\)-subgroup of \(G\) is abelian}
\quad\Longleftrightarrow\quad
\exists N\geq 1\ \forall k\geq N\ \forall g,h\in G,\qquad a_{p,k}(g,h)=1_G.
\]
Here \(N\) may depend on \(p\) and \(G\). Thus ``almost all'' means all but finitely many indices.

\textbf{Problem (R. Brandl).} For each prime \(p>2\), give an explicit sequence \((a_{p,k})_{k\geq 1}\) with the displayed property. The source reports an explicit sequence for \(p=2\). It does not specify the required meaning of ``explicit'' or whether a single algorithm taking \(p\) as input is intended. Its existence assertion and the reported \(p=2\) result are retained as supplied background, without independent verification.
\subsection{Short English statement}
For each odd prime \(p\), explicitly describe a sequence of two-variable group words such that a finite group has abelian Sylow \(p\)-subgroups exactly when every sufficiently late word is an identity in that group. The supplied record labels the problem open and separately reports an unverified AI claim of a computable construction; that claim is not treated here as an established solution.
\subsection{Sources}
\begin{itemize}
\item Local UnsolvedMath registry: permanent ID \(700708\), exact source problem code \texttt{catalogue entry}, title \emph{Kourovka Notebook Problem 11.15}; upstream revision \texttt{d371ee917ac9647fbb06cd6fc529deec6f3dda82}.
\item R. Brandl, proposer, 1990. The supplied record attributes the problem to \emph{The Kourovka Notebook (Unsolved Problems in Group Theory)}, 11th issue (1990), Problem 11.15, as included in the 21st issue edition (2026), arXiv:1401.0300v48. The registry identifies the notebook editors as E. I. Khukhro and V. D. Mazurov.
\item R. Brandl, \emph{J. Austral. Math. Soc.} \textbf{31} (1981), 464--469. The source cites this reference for an explicit sequence when \(p=2\); no article title is supplied.
\item Achyuth Jayadevan, \emph{Preprint} (2026), \url{https://zenodo.org/records/22852259}. The supplied background and status review explicitly label the reported computable construction an ``unverified AI solution'' claim. The status review says that the October 2026 notebook edition marks it with \(\#\).
\item Supplied status review dated \texttt{2026-10-16}: canonical status \texttt{open}, evidence category \texttt{unrefereed\_claim}. This is supplied dated curation, not an independently verified literature assessment.
\end{itemize}
\end{document}
